% Appendix B: certificates
\section{The computer-assisted parts}\label{app:cert}

The scripts are in the directory \texttt{certificates/} of the repository accompanying this paper. They are written in Python~3;
the interval computations use ball arithmetic from \texttt{python-flint} (Arb). Each subsection states what is checked, the
method, the command, and the expected output. Commands are run from the repository root unless stated otherwise.

Only two of them are steps of the proof: the Beurling--Saias certificate (Section~\ref{app:cert-bs}) and the small range of the
$X_0$-band bound (Section~\ref{app:cert-x0band}). The others are independent numerical cross-checks of statements proved in the
text (Theorem~\ref{thm:caplw}, Lemma~\ref{lem:sc}, Theorem~\ref{thm:svl}) or of the bookkeeping of Section~\ref{sec:logic}.

\subsection{The Beurling--Saias certificate (\texttt{bs-cert})}\label{app:cert-bs}
\emph{What it certifies.} The proof of Theorem~\ref{thm:bs} (and of Lemma~\ref{lem:bs-profile}) reduces the induction step to an
explicit inequality $\Sigma(\ell,u)\ge0$, where $\ell=\log Y$, $u=\ell/\log z$ and $\tau=\log t\ge10^8$. Two parts of this reduction
are computer-assisted.
\begin{itemize}
\item \emph{The region $30\tau\le\ell<10^{18}$, $u\ge2$.} A closed-form lower bound $S(\ell,u)\le\Sigma(\ell,u)$ is built from the
majorants of the proof and from bounds $r_{\mathrm{low}}\le r\le r_{\mathrm{up}}$ for the Dickman ratio $r=-\rho'/\rho$. The script
checks $S>0$ on cells covering the whole region, uniformly in $\tau\ge10^8$.
\item \emph{One-variable inequalities:} the positivity $c(u)\ge K'u^2\log(1+u)$ for $27\le u\le4\cdot10^5$, used for
$\ell\ge10^{18}$; the one-step inequality of Lemma~\ref{lem:smooth-v} for $1\le u\le10^6$; and that one-step inequality on
$u\in[1,2]$ including both extra terms of the proof: the excess $X(u)$ of the small heights $1\le Y'<72$ (present for
$u<1+\log72/L_z$), and a conservative allowance $2K_Mr(u)(e^{0.141}-1)$. Beyond these ranges they are proved analytically.
\end{itemize}
\emph{Method.}
\begin{itemize}
\item All arithmetic is ball arithmetic (Arb, through \texttt{python-flint} 0.9.0; module \texttt{ivarb.py}). Every elementary
operation and transcendental function is enclosed. The bounds $r_{\mathrm{low}}$ (from the recursion
$r(v)\ge r(v-1)/(1-e^{-r(v-1)})$, $r(2)=1/(2-2\log2)$) and $r_{\mathrm{up}}$ (from $r(x)\le\log((x+1)/\delta)/(1-\delta)$, optimised
over $\delta$) are computed in balls. Conversions to floating point are rounded outward and are used only to select cells.
\item Each cell is the product of an $\ell$-interval and a $u$-interval (or an interval in $v:=u-\ell^{1/4}$). The bound $S$ is
evaluated on the whole box, so no Lipschitz or interpolation argument is needed.
\item The range $\ell\in[3\cdot10^9,10^{18}]$ is split into 403 bands of ratio $1.05$. In each band, $u$-cells cover
$[2,\max(40,u_0/2)]$ and $v$-cells cover the rest up to the largest junction $\bar U(\ell)=U(\ell;\tau=10^8)$. The endpoints of the
$v$-cells are rounded outward, so the seams between $u$-cells and $v$-cells, and the top edge, are covered exactly for every $\ell$ in
the band. Cells that fail are bisected adaptively.
\item Every cell also checks $r(u-1)>\gamma_1$ and $r(u-1)>\gamma_{\max}$, which is the condition for the integrals of the second
majorant to converge.
\end{itemize}
\emph{How to run} (\texttt{pip install python-flint==0.9.0}):
\begin{verbatim}
python3 -I certificates/bs-cert/cert_run_arb_patched.py 1.05 10   # main region, 10 workers, ~20 s
python3 -I certificates/bs-cert/aux_cert_arb.py                   # the one-variable checks, ~2 s
\end{verbatim}
\emph{Expected output} (abridged; the first line is wrapped):
\begin{verbatim}
[patched seams] q=1.05 bands=403 cells=302013 failed=0 min certified l*S=2.491e+06
    at (82792994654.45967, 86932644387.18266, 200.1275416813517, 210.13391876541928, False)
(H'') c_A >= K'u^2 log(1+u): [27,400000] cells=9609 failed=0 min certified margin=0.189
Lemma V one-step (L*Q_V, L = max(100,u)): [1,1e+06] cells=13823 failed=0 min certified margin=1.724
Lemma V one-step on [1,2] with both extra terms: [1,2] cells=1387 failed=0 min certified margin=0.829
\end{verbatim}
That is: $302\,013$ cells (adaptive bisections included), $0$ failures, minimum certified $\ell\cdot S=2.491\cdot10^6$, at
$\ell\approx8.3\cdot10^{10}$, $u\in[200,210]$; for $c(u)$, $9\,609$ cells on $[27,4\cdot10^5]$, $0$ failures, minimum margin $0.189$;
for the one-step inequality of Lemma~\ref{lem:smooth-v}, $13\,823$ cells on $u\in[1,10^6]$, $0$ failures, minimum margin $1.724$;
and for that inequality on $u\in[1,2]$ with both extra terms, $1\,387$ cells, $0$ failures, minimum margin $0.829$. The script \texttt{cert\_run\_arb.py 1.05 10} gives the same
result. The double-precision twins \texttt{bs\_cert.py}, \texttt{cert\_run.py} and \texttt{aux\_cert.py} (minimum $2.598\cdot10^6$ for
the main region) are a cross-check only and are not part of the certificate. In the code, the first majorant is \texttt{betaI},
\texttt{alpha}; the second is \texttt{betaII}, \texttt{gam1}, \texttt{gmax}, \texttt{tailII}; the error terms are \texttt{E1},
\texttt{E2}, \texttt{E3} with \texttt{eta}; the tail is \texttt{ET}, \texttt{T}; and the $v$-cells are \texttt{vmode}.

\subsection{The jump at $X_0$, small range (\texttt{x0-band})}\label{app:cert-x0band}
\emph{What it certifies.} The small range of Lemma~\ref{lem:xb}: the bound $B(n)\le5/12$ for every integer $2\le n\le10^6$, where
\[
  B(n)=\frac1n\Bigl[\tfrac12\mu^2(n)+\tfrac12\sum_{d\mid n,\,1<d<n}\Bigl(2\mu^2(d)+\tfrac12\lambda\frac{X(d)}d\Bigr)\Bigr],\qquad
  X(d)=\sum_{e\mid d,\,1<e<d}\mu^2(e)\,e\,H_{\mathrm{ord}}(d/e),\qquad\lambda=1.0001.
\]
The maximum is attained only at $n=6$, where $B(6)=5/12$ exactly. The script also prints the suprema over $n\ge10^a$ for
$a=1,\dots,5$, and the analytic tail bound $(\tfrac12+1.64\sqrt n)/n+\tfrac14\lambda A(1.8)2^{0.2}n^{-0.2}$ at $n=10^6$, which is
$\le0.138$. For $n>10^6$ the lemma uses that analytic bound, not the computation.

\emph{Method.} Sieves compute, exactly in integers: $\mu^2$; $H_{\mathrm{ord}}(m)$ by $H_{\mathrm{ord}}(1)=1$,
$H_{\mathrm{ord}}(m)=\sum_{d\mid m,\,d<m}H_{\mathrm{ord}}(d)$; $X(d)$; and the divisor sums. The integer arrays are checked against
64-bit overflow (the maximum of $X$ is about $2.9\cdot10^8$). The only floating-point steps are the divisions $X(d)/d$ and sums of at
most $\tau(n)\le240$ nonnegative terms, whose relative rounding error is $<10^{-12}$; the script multiplies every $B(n)$ by
$1+10^{-9}$ before taking maxima, so every printed value is an upper bound. All terms are nonnegative and every integer quantity is
exact, so no interval arithmetic is needed.

\emph{How to run} (standard library only, about 6 seconds):
\begin{verbatim}
cd certificates/x0-band && python3 -I x0band_exact.py 1000000
\end{verbatim}
\emph{Expected output} (abridged):
\begin{verbatim}
N= 1000000  Hmax= 83934208  Xmax= 288745984  (int64 ok: True )
top: [(6, 0.41667), (12, 0.27778), (10, 0.25), (4, 0.25), (2, 0.25), (30, 0.23389), ...]
sup B(n), 10<=n<=N: 0.277781
sup B(n), 100<=n<=N: 0.088059
sup B(n), 1000<=n<=N: 0.022826
sup B(n), 10000<=n<=N: 0.007607
sup B(n), 100000<=n<=N: 0.002121
analytic tail bound for n>=N: 0.13738898318455534
\end{verbatim}
The companion script \texttt{x0band.py} is an earlier floating-point exploration up to $2\cdot10^4$; it is not part of the certificate.

\subsection{The capacity theorem: numerical cross-check (\texttt{capl-cert})}\label{app:cert-capl}
Theorem~\ref{thm:caplw} is proved analytically. The scripts in \texttt{certificates/capl-cert/} re-evaluate, in double-precision
floating point, the precise criteria that Steps~2--5 of that proof bound analytically: the degree cap $D^*$ computed from
$\binom{m+k}{k}$; the exact $j=\max\{j:e(r_0(D+1))^{2j+1}<c\}$; the row-weight bounds of Lemma~\ref{lem:rowweights}(b),(d); and
condition~\eqref{eq:twelvetwo} of Lemma~\ref{lem:lemmaP} with the exact absorption threshold $y$. They are not interval-certified, and
nothing in the proof depends on them. They use only the Python standard library and each takes a few seconds.
\begin{itemize}
\item \texttt{capl\_model.py}: the shared model (route-H criterion, Dickman-row criterion, Lemma~\ref{lem:lemmaP} criterion), in log
space. It prints nothing.
\item \texttt{analytic\_capl.py}: computes $s_H=5.9745$ (the root of $\gamma_*(s_H)=\tfrac13-10^{-5}$); checks that $G(\Lambda,s)$
decreases in $\Lambda$ for $\Lambda$ from $10^8$ to $10^{40}$ at sample values of $s$; and evaluates the precise criteria with
$c=\max(\tfrac14H_{\mathrm{std}}^{1/2},X_0^{\gamma(s)})$ on the grid $\Lambda\in\{10^8,10^9,10^{10},10^{12},10^{16},10^{20}\}$,
$s\in[-\log\Lambda,30]$ at step $0.01$. Expected output: \texttt{s\_H = 5.9745}, ``decreasing'' on every sample line, and
$4842$, $5072$, $5302$, $5763$, $6684$, $7605$ points with $0$ failures ($35\,268$ in all).
\item \texttt{thm417\_check\_refereeB.py}: an independent re-implementation of Steps~1, 2, 3 and~5 at $200\,000$ random pairs with
$\log_{10}\Lambda$ uniform in $[8,30]$ and $s\in[-\log\Lambda,60]$. Expected output:
\texttt{Step1 violations: 0 route-H violations: 0 (min lhs/rhs 1.6938) Dickman violations: 0 (min ratio 1.1897) Step5 violations: 0}.
\item \texttt{cert\_capl.py}: an earlier cell-by-cell check at $\Lambda=10^8$ only, with a piecewise-linear ramp of maximum $5.95$ in
place of $\gamma$, covering $\ell\in[\log\Lambda,2\cdot10^4]$ with cells of width $0.002$ up to $\ell=60$ and $0.05$ beyond; each cell is
checked at its worst corners. Expected output: \texttt{cells 20790 on [18.421, 60]: failures 0; max slope dlogH/dL = 2.94e-06} and
\texttt{cells 398801 on [60, 2e4]: failures 0} ($419\,591$ cells in all). It is superseded by the analytic, $\Lambda$-uniform proof.
\item \texttt{needs.py}: a diagnostic that prints, at sample $\ell$, the least $\log c$ needed by the precise criteria and the largest
allowed by Lemma~\ref{lem:lemmaP}. It checks nothing.
\item \texttt{raise\_max.py}: the value $1.3\cdot10^{-8}$ quoted in Step~6 of the proof of Theorem~\ref{thm:caplw}. It computes the
raise $\max(0,2\log M_1-8\ell^2)/L$ of $H$ over $H_{\mathrm{std}}$, with $c$, $\gamma(s)$ and $M_1=r_0+\varphi(r_0)+1$ exactly as in that
theorem (in log space), on the grid $\log_{10}\Lambda\in[8,30]$ at step $0.05$ and $s\in[-\log\Lambda,30]$ at step $0.01$
($3\,252\,514$ points), and compares the maximum with the analytic bound $10^{-6}$, which is all the proof uses. Expected output
(about 10 seconds):
\begin{verbatim}
s_H = 5.9745; points = 3252514
max (log H - log H_std)/L = 1.292e-08 at log10 Lambda = 8.00, s = -9.92; analytic bound 1e-6: OK
\end{verbatim}
\end{itemize}


\subsection{Lemma~\ref{lem:sc}: numerical check (\texttt{lemma-sc})}\label{app:cert-sc}
Lemma~\ref{lem:sc} is proved in full in Appendix~\ref{app:rankfull}; these scripts are an independent corroboration, not a step of
the proof. They work in exact rational arithmetic (Python's \texttt{fractions}).
\begin{itemize}
\item \emph{What is tested.} Random rank-full integer exponent matrices $E$ ($r\le8$ rows, up to $r+6$ columns, entries in
$\{0,1,2,3\}$ at random densities; columns that vanish identically are removed, since primes absent from every row must not be
counted) whose greedy latest-first pivots form at least two stages. For every \emph{minimum} cover $C$ of the row supports,
\texttt{sc\_check.py} searches by brute force for nested row sets $R_{\le1}\subset R_{\le2}\subset\cdots$ with
$|R_{\le t}|=|P_{\le t}|$, $E_{R_{\le t},P_{\le t}}$ invertible, and every pivot of $P_{\le t}$ and every late free column of $F_t$ covered
by a row of $C\cap R_{\le t}$. \texttt{sc\_greedy.py} tests the explicit construction of the proof: the lexicographically first
bases of the row matroids $M_t$, with the rows of $C$ first, checked for nesting, invertibility and coverage.
\texttt{sc\_check\_refereeA.py} is an independent re-implementation over all \emph{minimal} covers; it also shows that the order
matters, since putting the rows of $C$ last fails coverage in $1238$ of $2342$ cases.
\item \emph{Commands and expected output} (each run prints one dictionary and takes seconds to a minute):
\begin{verbatim}
python3 -I certificates/lemma-sc/sc_check.py 1
  {'tested': 10615, 'single': 20820, 'all_ok': 10615, 'some_fail': 0, 'all_fail': 0}
python3 -I certificates/lemma-sc/sc_check.py 7 20000 7
  {'tested': 5585, 'single': 10441, 'all_ok': 5585, 'some_fail': 0, 'all_fail': 0}
python3 -I certificates/lemma-sc/sc_check.py 11 60000 5
  {'tested': 15933, 'single': 31583, 'all_ok': 15933, 'some_fail': 0, 'all_fail': 0}
python3 -I certificates/lemma-sc/sc_greedy.py 5 30000 8
  {'ok': 13433}
python3 -I certificates/lemma-sc/sc_check_refereeA.py
  configs 842 {'ok': 2342, 'fail': 0, 'advfail': 1238, 'nest': 0}
\end{verbatim}
Here \texttt{tested} counts the multi-stage configurations (single-stage ones are skipped and counted as \texttt{single}), and
\texttt{all\_ok} those for which the nested sets exist for every minimum cover; \texttt{ok} counts the (configuration, cover) pairs
for which the explicit construction works, and \texttt{advfail} the coverage failures of the reversed order.
\end{itemize}

\subsection{The staged volume lemma: structural check (\texttt{svl-stress})}\label{app:cert-svl}
Theorem~\ref{thm:svl} is proved in Appendix~\ref{app:rankfull}; this script is an independent sanity check of the structural facts
its proof uses, not a step of the proof. In exact rational arithmetic, on random and adversarial rank-full configurations with
$8$--$12$ rows and up to nine stages, it checks that the stage rows exist; that column masses are $\ge1$ on the stage prefix; that
late free primes are covered; that a late free column is determined by its stage-prefix entries, with vanishing coefficient in the
remaining rows; that the cumulative determinants are nonzero integers; and that the per-row block ranges are disjoint.
\begin{verbatim}
python3 -I certificates/svl-stress/svl_stress.py        # a few minutes, standard library only
  rand configs 60 max stages T 1 violations 0
  late configs 60 max stages T 9 violations 0
  refA configs 60 max stages T 2 violations 0
  nest configs 34 max stages T 4 violations 0
\end{verbatim}

\subsection{Non-circularity of the dependency graph (\texttt{depcheck})}\label{app:cert-dep}
Section~\ref{sec:logic} asserts that the dependencies between the results of this paper used at the \emph{same} time form an acyclic
graph, every loop closing only through an input used at a strictly earlier time or a strictly smaller height. The script
\texttt{certificates/depcheck/depcheck.py} encodes the full dependency graph, one edge per input, each edge marked ``same time'' or
``strictly earlier'' as in the table of Section~\ref{sec:logic} (its node names follow the numbering of the working draft, e.g.\
\texttt{T4.3} for Theorem~\ref{thm:queue}). It runs a depth-first search for a cycle among the same-time edges.
\begin{verbatim}
python3 -I certificates/depcheck/depcheck.py
  same-time subgraph acyclic: True | nodes 77 | edges 168 | strictly-earlier edges 27
\end{verbatim}
The time-ordering of the strictly-earlier edges is argued in Section~\ref{sec:noncirc}; the script checks only the bookkeeping.
